\documentclass[10pt,a4paper]{amsart}
\usepackage[margin=19mm]{geometry}
\usepackage{amsmath,amssymb,mathtools}
\usepackage[scr=rsfs,cal=euler]{mathalfa}
\usepackage{xfrac}
\usepackage[unicode,hidelinks,bookmarks=false]{hyperref}
\newcommand{\ie}{i.e.\ }
\newcommand{\cf}{cf.\ }
\newcommand{\resp}{respectively\ }
\newcommand{\Subsection}{\subsection}
\providecommand{\nobreakdash}{\nobreak\hbox{-}}
\setlength{\emergencystretch}{2em}
\multlinegap0pt
\usepackage{mathtools}
\usepackage{amssymb}
\usepackage[shortlabels]{enumitem}
\newtheorem{theorem}{Theorem}
\title{Zero modes on product Riemannian manifolds}
\author{Jurgen Julio-Batalla}
\date{Source: arXiv:2603.23242v1 (CC BY 4.0)}
\begin{document}
\maketitle

\noindent\emph{Source and licence.} Zero modes on product Riemannian manifolds. Version arXiv:2603.23242v1 (CC BY 4.0), distributed by arXiv under the Creative Commons Attribution 4.0 International licence, http://creativecommons.org/licenses/by/4.0/, as recorded in the arXiv licence metadata for this exact version and shown on the abstract page https://arxiv.org/abs/2603.23242v1. Full licence notice: \texttt{licenses/arxiv-2603.23242-CC-BY-4.0.txt}.\par\smallskip

\noindent\textit{Editorial scope.} Counted unit: the article's theorem theorem (source0.tex lines 303--306). The article's complete original proof of it is delivered with it (source0.tex lines 308--337, one \texttt{proof} environment of 30 lines). No mathematics is written, repaired or re-derived: the statement and the proof are the article's own lines, in the article's own notation, and no retained line is substituted or re-flowed.  Retained uncounted as the source-local context the proof needs: lines 291--293.  Nothing was omitted from any retained span, so the package carries no omission record.  No retained line contains a citation, so the wrapper carries no bibliography and no bibliography entry.  No retained line contains a figure environment, an image-inclusion command or drawing code, so no image is delivered and the package declares no asset dependency.  Editorial formatting only, in addition: an AMS \texttt{amsart} preamble that restates the article's own declarations and macros used by the retained lines, namely the environments \texttt{theorem} on separate counters, together with the standard packages those lines need.  The retained environments print in this document as Theorem 1 in order of appearance, whereas the article prints the same items with the numbering of its own full document.  The complete native source of the article as distributed in the arXiv e-print for this exact version is bundled unchanged as \texttt{sources/197066/source0.tex}.
This operator satisfies that \begin{equation}\label{decompositionT}
|T\varphi|^2=|\nabla\varphi|^2-\frac{1}{n_1}|D_{(1)}\varphi|^2-\frac{1}{n_2}|D_{(2)}\varphi|^2
\end{equation}

\begin{theorem}
Let $(M^n=M_1^{n_1}\times M_2^{n_2},g=g_1+g_2,\sigma)$ be a product of closed spin manifolds of dimensions $n_1\leq n_2$. Assume that the scalar curvature $s_g$ is positive. Let $(\varphi,A)$ be a non trivial solution of the zero mode equation $$D_g\varphi=i\lambda A\cdot_g\varphi,\quad\int_M |A|^ndv_g=1. $$
If $|\varphi|^2$ is non-increasing with respect to $|A|$ then $$\lambda^2\geq \frac{n_2}{4(n_2-1)}Y(M^n,[g]).$$
\end{theorem}

\begin{proof}
By the Schrödinger-Lichnerowitz formula $$0=\int_M\left(- \left\langle D^2_g\varphi,\varphi\right\rangle+\frac{s_g}{4}|\varphi|^2+|\nabla \varphi|^2\right)dv_g$$
Using Stoke Theorem and the zero mode equation we have $$0=\int_M\left(- \lambda^2|A|^2|\varphi|^2+\frac{s_g}{4}|\varphi|^2+|\nabla \varphi|^2\right)dv_g.$$
Taking the Penrose-type operator $T$ from the section 2, we decompose the term $|\nabla\varphi|^2$ as \eqref{decompositionT} in order to get
$$0=\int_M\left( - 4\lambda^2|A|^2+s_g\right)|\varphi|^2+\int_M 4|T\varphi|^2+\int_M\left(\frac{4}{n_1}|D_{(1)}\varphi|^2+\frac{4}{n_2}|D_{(2)}\varphi|^2\right).$$
Since each $D_{(i)}$ is self adjoint and $D^2=D^2_{(1)}+D^2_{(2)}$ we see $$\int_M|D_{(2)}\varphi|^2=\int_M|D\varphi|^2-\int_M|D_{(1)}\varphi|^2.$$
Therefore, from the two last equations we obtain
 
$$0=\int_M |\varphi|^2\left(4\left(\frac{1-n_2}{n_2}\right)\lambda^2|A|^2+s_g\right)+\int_M 4|T\varphi|^2+\int_M 4\left( \frac{1}{n_1}-\frac{1}{n_2}\right)|D_{(1)}\varphi|^2.$$
On the other hand, for any positive function $u$ we have
$$\int_Ma_n\left\langle\nabla ln(u),\nabla |\varphi|^2\right\rangle=\int_M-a_n|\varphi|^2\left(\frac{\Delta u}{u}-\frac{|\nabla u|^2}{u^2}\right).$$

Adding the last two equations we get the following integral identity (*)
\begin{align*}
0=\int_M\frac{|\varphi|^2}{u}\left(-a_n\Delta u+s_gu+4\left(\frac{1-n_2}{n_2}\right)\lambda^2|A|^2u\right)+\int_M 4|T\varphi|^2+\int_M 4\left( \frac{1}{n_1}-\frac{1}{n_2}\right)|D_{(1)}\varphi|^2\\
+\int_M a_n|\varphi|^2\frac{|\nabla u|^2}{u^2}-\int_Ma_n\left\langle\nabla ln(u),\nabla |\varphi|^2\right\rangle.
\end{align*}
Now consider the constant $I(M,g,|A|^2)$ defined by $$I(M,g,|A|^2)=\inf_{H^1(M^n)}\frac{\int_M\left( a_n|\nabla u|^2+s_gu^2\right)dv_g}{\int_M (|A|u)^2dv_g}.$$
i.e. the first eigenvalue of the linear weighted problem $-a_n\Delta u+s_gu=I|A|^2u.$

Let $u_1$ be the positive smooth function which realizes $I(M,g,|A|^2)$. By the non increasing assumption  on $|\varphi|^2$, we have $$-\int_Ma_n\left\langle\nabla ln(u_1),\nabla |\varphi|^2\right\rangle\geq 0.$$
In particular, inserting $u_1$ in the integral identity (*) we conclude that $$4\left(\frac{n_2-1}{n_2}\right)\lambda^2< I(M,g,|A|^2)$$
implies that $\varphi$ must be trivial.
Hence  $$4\left(\frac{n_2-1}{n_2}\right)\lambda^2\geq I(M,g,|A|^2)$$
To obtain a lower bound for $I$ note from normalization of $|A|$ and Hölder's inequality that $$\int_M (|A|u)^2dv_g\leq \left(\int_Mu^{2n/(n-2)}dv_g\right)^{(n-2)/n},$$
and hence

 $$I\geq \inf_{H^1(M)} \frac{\int_M\left( a_n|\nabla u|^2+s_gu^2\right)dv_g}{\left(\int_Mu^{2n/(n-2)}dv_g\right)^{(n-2)/n}}=Y(M^n,[g]). $$ 
Therefore $$\lambda^2\geq \frac{n_2}{4(n_2-1)}Y(M^n,[g]).$$
\end{proof}

\end{document}
